% -------------------------------------------------------------------------

\subsection{Теоретический вопрос IV.1: Выведите SVD из произведения матриц $A^\top A$}
\textbf{Формулировка вопроса:} Конструктивно выведите сингулярное разложение (SVD) произвольной вещественной матрицы $A
\in \mathbb{R}^{m \times n}$ через спектральный анализ матрицы $A^\top A$.

\textbf{Ответ и строгий вывод:}

Пусть дана произвольная вещественная матрица $A \in \mathbb{R}^{m \times n}$ ранга $\rank(A) = r \le \min(m, n)$.
Рассмотрим квадратную матрицу $S \coloneqq A^\top A \in \mathbb{R}^{n \times n}$.
\begin{enumerate}
    \item \textbf{Свойства матрицы $A^\top A$:}
    \begin{itemize}
        \item Симметричность: $(A^\top A)^\top = A^\top A$.
        \item Неотрицательная определенность: для любого $x \in \mathbb{R}^n$ имеем $x^\top (A^\top A) x =
        (Ax)^\top (Ax) = \|Ax\|_2^2 \ge 0$.
    \end{itemize}
    \item По вещественной спектральной теореме матрица $A^\top A$ обладает ортонормированным базисом из
    собственных векторов $(v_1, v_2, \dots, v_n)$ в $\mathbb{R}^n$, а все ее собственные значения вещественны и
    неотрицательны.
    Упорядочим их по невозрастанию:
    \[ \lambda_1 \ge \lambda_2 \ge \dots \ge \lambda_r > \lambda_{r+1} = \dots = \lambda_n = 0. \]
    При этом:
    \[ A^\top A v_i = \lambda_i v_i, \quad v_i^\top v_j = \delta_{ij}. \]
    Векторы $v_1, \dots, v_n$ образуют столбцы ортогональной матрицы $V \in \mathbb{R}^{n \times n}$ (
    $V^\top V = I_n$).
    \item \textbf{Определение сингулярных чисел:} Положим
    \[ \sigma_i \coloneqq \sqrt{\lambda_i} > 0 \quad \text{для } i = 1, \dots, r. \]
    \item \textbf{Построение левых сингулярных векторов $u_i$:} Для каждого $i \in \{1,\dots, r\}$ определим вектор:
    \[ u_i \coloneqq \frac{1}{\sigma_i} A v_i \in \mathbb{R}^m. \]
    Проверим скалярные произведения векторов $u_i$ и $u_j$:
    \[ u_i^\top u_j = \left(\frac{1}{\sigma_i} A v_i\right)^\top \left(\frac{1}{\sigma_j} A v_j\right) =
    \frac{1}{\sigma_i \sigma_j} v_i^\top (A^\top A v_j) = \frac{1}{\sigma_i \sigma_j} v_i^\top (\lambda_j v_j) =
    \frac{\lambda_j}{\sigma_i \sigma_j} (v_i^\top v_j). \]
    Так как $v_i^\top v_j = \delta_{ij}$ и $\lambda_i = \sigma_i^2$:
    \[ u_i^\top u_j = \frac{\sigma_j^2}{\sigma_i \sigma_j} \delta_{ij} = \delta_{ij}. \]
    Следовательно, система векторов $\{u_1, \dots, u_r\}$ строго ортонормирована в $\mathbb{R}^m$.
    \item Если $r < m$, дополним $\{u_1, \dots, u_r\}$ процедурой Грама-Шмидта до полного ортонормированного базиса
    $\{u_1, \dots, u_m\}$ пространства $\mathbb{R}^m$.
    \item Составим из них ортогональную матрицу $U = [u_1 \dots u_m]
    \in \mathbb{R}^{m \times m}$ ($U^\top U = I_m$).
    \item По построению:
    \[ A v_i = \sigma_i u_i \quad (1 \le i \le r), \quad A v_i = 0 \quad (r < i \le n). \]
    В матричной форме это записывается как:
    \[ A V = U \Sigma \iff A = U \Sigma V^\top, \] что полностью завершает конструктивный вывод SVD\@.
\end{enumerate}

% -------------------------------------------------------------------------

\subsection{Теоретический вопрос IV.2: Почему $A^\top A$ и $A A^\top$ имеют общие ненулевые eigenvalues?}
\textbf{Формулировка вопроса:} Докажите, что матрицы $A^\top A$ и $A A^\top$ имеют в точности одинаковые ненулевые
собственные значения с учётом кратностей.

\textbf{Ответ и строгое доказательство:}

Пусть $A \in \mathbb{R}^{m \times n}$.
Матрица $A^\top A$ имеет размер $n \times n$, а матрица $A A^\top$ имеет размер $m \times m$.
\begin{enumerate}
    \item Пусть $\lambda \ne 0$ — ненулевое собственное значение матрицы $A^\top A$, и $v \in \mathbb{R}^n$ —
    соответствующий ненулевой собственный вектор:
    \[ (A^\top A) v = \lambda v, \quad v \ne 0. \]
    \item Умножим обе части равенства слева на матрицу $A$:
    \[ A (A^\top A v) = A (\lambda v) \iff (A A^\top) (A v) = \lambda (A v). \]
    \item Определим вектор $w \coloneqq A v \in \mathbb{R}^m$.
    Проверим, отличен ли он от нуля:
    Вычислим квадрат евклидовой нормы вектора $w$:
    \[ \|w\|_2^2 = w^\top w = (A v)^\top (A v) = v^\top (A^\top A v) = v^\top (\lambda v) = \lambda \|v\|_2^2. \]
    Так как $\lambda > 0$ и $v \ne 0$, имеем $\|w\|_2^2 > 0$, откуда $w \ne 0$.
    \item Следовательно, вектор $w = Av$ является \textbf{ненулевым собственным вектором} матрицы $A A^\top$,
    отвечающим тому же самому собственному значению $\lambda$.
    \item Проводя аналогичное рассуждение в обратную сторону: если $(A A^\top) u = \mu u$ при $\mu \ne 0, u \ne 0$,
    то умножение слева на $A^\top$ даёт $(A^\top A)(A^\top u) = \mu (A^\top u)$ с ненулевым вектором $A^\top u$.
    \item \textbf{Совпадение геометрических кратностей:} Отображение $v \mapsto Av$ устанавливает линейный
    изоморфизм между собственным подпространством матрицы $A^\top A$ для значения $\lambda$ и собственным
    подпространством матрицы $A A^\top$ для того же $\lambda$.
    Следовательно, их размерности совпадают.
\end{enumerate}

% -------------------------------------------------------------------------

\subsection{Теоретический вопрос IV.3: Объясните outer-product SVD form}
\textbf{Формулировка вопроса:} Объясните представление SVD в виде суммы внешних произведений (outer-product form), ее
блочную алгебру и физический смысл.

\textbf{Ответ:}

Используя редуцированное сингулярное разложение $A = U_r \Sigma_r V_r^\top$, где $U_r = [u_1 \dots u_r]$, $V_r = [v_1
\dots v_r]$ и $\Sigma_r = \diag(\sigma_1, \dots, \sigma_r)$, распишем матричное произведение по правилу блочного
умножения «столбец на строку»:
\[ A = U_r \Sigma_r V_r^\top = [u_1 \dots u_r] \begin{pmatrix}
                                                           \sigma_1 v_1^\top \\ \vdots \\ \sigma_r v_r^\top
\end{pmatrix} = \sum_{i=1}^r \sigma_i u_i v_i^\top. \]

\textbf{Анализ элементарного слагаемого $u_i v_i^\top$:}
\begin{enumerate}
    \item $u_i \in \mathbb{R}^m$ ($m \times 1$), $v_i^\top \in \mathbb{R}^{1 \times n}$ ($1 \times n$).
    \item Произведение $u_i v_i^\top$ является матрицей размера $m \times n$ и имеет
    \textbf{ранг ровно 1} (столбцы пропорциональны $u_i$, строки пропорциональны $v_i^\top$).
    \item Матрицы $u_i v_i^\top$ взаимно ортогональны относительно скалярного произведения Фробениуса $\langle A, B
    \rangle_F = \tr(A^\top B)$:
    \[ \tr\left((u_i v_i^\top)^\top (u_j v_j^\top)\right) = \tr(v_i u_i^\top u_j v_j^\top) = (u_i^\top u_j)
    \tr(v_i v_j^\top) = \delta_{ij} (v_j^\top v_i) = \delta_{ij}. \]
\end{enumerate}

\textbf{Физическая и концептуальная интерпретация:} Внешняя форма SVD показывает, что любое сложное преобразование $A$
состоит из параллельной суммы $r$ элементарных независимых каналов:
\begin{itemize}
    \item $v_i^\top x$ — измерение проекции входа на входной базисный вектор $v_i$;
    \item $\sigma_i$ — коэффициент усиления (амплитуда, сила сигнала в данном канале);
    \item $u_i$ — синтез выходного направления в пространстве отклика.
\end{itemize}

\newpage

% =========================================================================
% ЧАСТЬ 2: ПРАКТИЧЕСКИЕ ВОПРОСЫ (PRACTICE)
% =====================================================================